% from ega4-18-fr.tex
\oldpage[IV-4]{153}

puisque le point fermé de $\Spec(A'')$ est le seul point au-dessus du point fermé de $\Spec(A)$. Or, lorsque $A$ est géométriquement unibranche, $k''$ est une extension radicielle de $k(A)$, donc $\Spec(k(B)\otimes_{k(A)}k'')$ est radiciel sur $\Spec(k(B))$ et ne comprend donc qu'un seul point. Lorsque $A$ est unibranche et $k(B)$ extension primaire de $k(A)$, $\Spec(k(B)\otimes_{k(A)}k'')$ est irréductible (4.3.2) et comme c'est un espace fini et discret (I, 6.4.4) il est encore réduit à un point, ce qui montre que dans les deux cas envisagés $B''$ est local. Si $Y''=\Spec(A'')$, $X''=\Spec(B'')=X\times_Y Y''$, le morphisme $f''=f_{(Y'')}:X''\to Y''$ est plat et a ses fibres géométriquement réduites; en outre, si $\eta''$ est le point générique de $Y''$, $k(\eta'')=k(\eta)=K$ par définition, donc les fibres $X_\eta$ et $X''_{\eta''}$ sont isomorphes. Il suffit donc de prouver que $X''_{\eta''}$ est connexe.

\begin{remark}[18.9.7.1]
\label{IV.18.9.7.1}
Lorsque l'anneau $A$ est \emph{japonais}, on peut, dans ce qui précède, prendre $A''=A'$ par définition (0, 23.1.1); on voit donc que dans ce cas on serait même ramené à prouver le théorème lorsque $A$ est \emph{intégralement clos}.
\end{remark}

\begin{env}[18.9.7.2]
\label{IV.18.9.7.2}
Supposons donc désormais que $Y$ soit intègre et géométriquement unibranche. Notons que si $T$ est une partie fermée de $X$ telle que $f(T)$ soit rare dans $Y$, $T$ ne contient aucun point maximal de $X$ puisque $f$ est plat (2.3.4), donc est rare dans $X$; comme $X$ est un schéma local, donc connexe, pour prouver que $X-T$ est connexe, il suffit, en vertu de (15.5.6.1) de montrer que pour tout $x\in X$ tel que $f(x)\neq\eta$, $\Spec(\sh{O}_{X,x})-\{x\}$ est connexe. Posons $y=f(x)$, $A_1=\sh{O}_{Y,y}$, $B_1=\sh{O}_{X,x}$, $Y_1=\Spec(A_1)$, $X_1=\Spec(B_1)$, $f_1:X_1\to Y_1$ le morphisme correspondant à l'homomorphisme local $A_1\to B_1$ déduit de $\rho$; il est clair que $f_1$ est plat, et il résulte de (4.6.1) que ses fibres sont géométriquement réduites; en outre $A_1$ est intègre et géométriquement unibranche (6.15.1) et $\dim(A_1)\geq 1$. On est ainsi ramené à prouver le
\end{env}

\begin{lemma}[18.9.7.3]
\label{IV.18.9.7.3}
Soient $A$, $B$ deux anneaux locaux noethériens, $Y=\Spec(A)$, $X=\Spec(B)$, $\rho:A\to B$ un homomorphisme local, $f:X\to Y$ le morphisme correspondant. On suppose vérifiées les conditions suivantes:
\begin{enumerate}
  \item[(i)] $f$ est plat et toutes les fibres $X_y$ ($y\in Y$) sont géométriquement réduites.
  \item[(ii)] $A$ est intègre, géométriquement unibranche et $\dim(A)\geq 1$.
\end{enumerate}
Alors, si $b$ est le point fermé de $X$, $X-\{b\}$ est connexe.
\end{lemma}

Notons d'abord qu'en vertu du théorème de Hartshorne (5.10.7), $X-\{b\}$ est connexe si l'on a $\operatorname{prof}(B)\geq 2$. Or on a (6.3.1)
\begin{equation}
  \operatorname{prof}(B)=\operatorname{prof}(A)+\operatorname{prof}(B\otimes_A k)
  \tag{18.9.7.4}
  \label{IV.18.9.7.4}
\end{equation}
en désignant par $k$ le corps résiduel de $A$. D'autre part, puisque $A$ est intègre et $\dim(A)\geq 1$, on a $\operatorname{prof}(A)\geq 1$, donc $\operatorname{prof}(B)\geq 2$ \emph{sauf} lorsque $\operatorname{prof}(B\otimes_A k)=0$; d'ailleurs $B\otimes_A k$ est réduit par (i), donc la relation $\operatorname{prof}(B\otimes_A k)=0$ signifie que $B\otimes_A k$ est un \emph{corps} (0, 16.4.7). On supposera donc désormais que cette dernière condition est vérifiée. On commencera par traiter un cas où la démonstration est très simple.

A) \emph{Cas où $A$ est intégralement clos.} Alors, si $\dim(A)\geq 2$, on a $\operatorname{prof}(A)\geq 2$ (0, 16.5.1), donc $\operatorname{prof}(B)\geq 2$. On n'a donc qu'à considérer le cas où $\dim(A)=1$ et où $B\otimes_A k$ est un corps; $A$ est alors un anneau de valuation discrète, donc régulier, et puisque $B\otimes_A k$

